\chapter{Órfãos --- Administração de BD e Temas Coringa}
\label{cap:orfaos}

\begin{conceito}[O que é este capítulo]
Bloco ``coringa'' do quiz: o que não encaixa nos outros capítulos. Na
prática, a maior parte das questões daqui é \textbf{administração de banco de
dados (DBA)} e \textbf{administração de dados} --- conteúdo do edital
(``noções de administração de dados e de banco de dados; backup, restauração;
otimização de performance'') que o capítulo de Banco de Dados, focado em
modelagem/SQL, não detalha. O resto são temas avulsos: arquitetura de
computadores, informática, sistemas de informação, IA/ML e blockchain. Nas
provas recentes a FGV vem carregando em IA/ML e cloud --- trate como tema
quente.
\end{conceito}

\section{Administração de Dados $\times$ Administração de Banco de Dados}

\begin{center}
\begin{tabular}{@{}p{2.2cm}p{4.6cm}p{4.6cm}@{}}
\toprule
& \textbf{Administração de Dados (AD)} & \textbf{Administração de BD (DBA)} \\
\midrule
Foco & \textbf{negócio/lógico}: o dado como recurso corporativo & \textbf{técnico/físico}: o SGBD funcionando \\
Faz & modelagem conceitual, políticas, governança, dicionário de dados & instalação, tuning, backup, segurança, disponibilidade \\
Papel & estratégico, independe de SGBD & operacional, ligado ao produto (Oracle, etc.) \\
\bottomrule
\end{tabular}
\end{center}

\begin{pegadinha}
AD é conceitual/negócio; DBA é físico/operacional. A FGV troca os dois papéis
entre si.
\end{pegadinha}

\section{Recuperação e transações (ARIES)}

\textbf{ARIES} (protocolo de recuperação da maioria dos SGBDs) usa
\textbf{WAL} (Write-Ahead Logging --- o log vai para disco \textbf{antes} do
dado) e três fases na recuperação após falha, \textbf{nesta ordem}:

\begin{enumerate}
\item \textbf{Analysis} (Análise): reconstrói o estado a partir do último
checkpoint.
\item \textbf{Redo} (Refazer): reaplica \textbf{todas} as operações
registradas (até as de transações que não commitaram) --- ``repeating
history''.
\item \textbf{Undo} (Desfazer): desfaz as transações que não commitaram.
\end{enumerate}

WAL garante a durabilidade (D do ACID) sem gravar o dado imediatamente.

\begin{pegadinha}
Ordem Analysis $\to$ Redo $\to$ Undo --- a FGV inverte Redo/Undo.
\end{pegadinha}

\section{Armazenamento físico e desempenho}

\begin{itemize}
\item \textbf{Tablespace} (Oracle) / filegroup: unidade lógica de
armazenamento que agrupa objetos em arquivos físicos. Separa dados, índices,
temporário.
\item \textbf{Índices} (B+ Tree, bitmap, hash): aceleram busca ao custo de
escrita. Bitmap é bom para \textbf{baixa cardinalidade} (poucos valores
distintos); B+ Tree para alta cardinalidade e faixas.
\item \textbf{Otimizador de consultas (query optimizer):} escolhe o
\textbf{plano de execução} (ordem de joins, uso de índice). Planos ruins
geralmente vêm de \textbf{estatísticas desatualizadas} --- atualizar
estatísticas costuma resolver.
\item \textbf{Particionamento} (horizontal/sharding, vertical): divide
tabela grande para escalar e melhorar desempenho.
\end{itemize}

\section{Backup e recuperação}

\begin{center}
\begin{tabular}{@{}p{2.6cm}p{8.4cm}@{}}
\toprule
\textbf{Tipo} & \textbf{O que copia} \\
\midrule
Completo (full) & tudo \\
Incremental & só o que mudou \textbf{desde o último backup (de qualquer tipo)} \\
Diferencial & tudo que mudou \textbf{desde o último backup completo} \\
\bottomrule
\end{tabular}
\end{center}

\textbf{RPO} (Recovery Point Objective): quanto de dado se aceita perder
(janela). \textbf{RTO} (Recovery Time Objective): em quanto tempo o serviço
volta. Backup \textbf{quente} (online, banco no ar) $\times$ \textbf{frio}
(offline).

\begin{pegadinha}
Incremental (desde o último \textbf{qualquer}) $\times$ diferencial (desde o
último \textbf{completo}) --- a diferencial cresce até o próximo full. A FGV
troca essa definição com frequência.
\end{pegadinha}

\section{Segurança e acesso no SGBD}

\textbf{GRANT/REVOKE} (DCL), \textbf{roles/papéis} (RBAC), \textbf{views}
para restringir colunas, \textit{row-level security}, auditoria. Princípio do
\textbf{menor privilégio}: por exemplo, o gerente de RH acessa só o que
precisa.

\section{Temas coringa genuínos}

\begin{itemize}
\item \textbf{Arquitetura de computadores:} CPU (Unidade de Controle
$\times$ ULA), ciclo de instrução (busca-decodifica-executa), pipeline,
memória (registrador $\to$ cache $\to$ RAM $\to$ disco, mais rápida e cara no
topo), sistemas de numeração (binário, octal, hexadecimal --- conversão dígito
a dígito).
\item \textbf{Noções de informática:} pacote office, sistemas de arquivos,
atalhos.
\item \textbf{Sistemas de informação:} SIG, SPT (transacional), SAD/SSD
(decisão), ERP, CRM, SCM --- o nível gerencial que cada um atende.
\end{itemize}

\section{Onde cada tema coringa está detalhado}

Este é o capítulo mais transversal do material: como o quiz aponta o capítulo
pela \textbf{tag} da questão, tudo que é \texttt{orfaos} cai aqui --- mesmo
quando o conteúdo mora, com razão, em outro capítulo. Use este mapa em vez de
procurar de página em página.

\begin{center}
\begin{tabular}{@{}p{4.4cm}p{6.6cm}@{}}
\toprule
\textbf{Tema} & \textbf{Onde está} \\
\midrule
Blockchain (hash do bloco anterior, raiz de Merkle, \emph{nonce},
imutabilidade encadeada) & Capítulo \ref{cap:programacao} \\
Servidor web $\times$ servidor de aplicação & Capítulo \ref{cap:arquitetura},
primeira seção \\
2PC (two-phase commit) e o contraste com Raft & Capítulo
\ref{cap:arquitetura} \\
Virtualização: hipervisor Tipo 1 $\times$ Tipo 2 $\times$ contêiner &
Capítulo \ref{cap:arquitetura} \\
Sistema operacional: paginação/memória virtual, E/S bloqueante, troca de
contexto & Capítulo \ref{cap:arquitetura} \\
Low-code / no-code & Capítulo \ref{cap:programacao} \\
RPA (Robotic Process Automation) & Capítulos \ref{cap:eng-software} e
\ref{cap:programacao} \\
Os Vs do Big Data & Capítulo \ref{cap:banco-dados} \\
Internet $\times$ intranet $\times$ extranet $\times$ portal & Capítulo
\ref{cap:redes} (e a pegadinha neste capítulo) \\
IA generativa, LLM, RAG, ética e regulação & Capítulo \ref{cap:atualidades} \\
\bottomrule
\end{tabular}
\end{center}

Três desses valem uma frase aqui, porque são exatamente o que a questão
cobra e o candidato costuma resolver no reflexo errado:

\begin{itemize}
\item \textbf{2PC:} um \textbf{coordenador} pergunta a todos os participantes
se podem confirmar (fase \emph{prepare}) e só ordena o \emph{commit} com
\textbf{unanimidade} --- um único ``não'' aborta a transação inteira. Não é
maioria simples (isso é consenso, tipo Raft), nem confirmação independente de
cada nó.
\item \textbf{RPA:} software que \textbf{imita o usuário} (clique, digitação,
leitura de tela) em tarefas digitais repetitivas e \textbf{baseadas em
regras}. Não é IA, não é robô físico, não substitui banco de dados.
\item \textbf{Low-code:} desenvolvimento por \textbf{modelagem visual e
configuração}, com o \textbf{mínimo} de código manual --- \emph{no-code}
elimina o código de vez. Nenhum dos dois ``dispensa desenvolvedor
profissional'': esse absoluto é o distrator.
\end{itemize}

\section{IA/ML (tema quente neste bloco)}

\begin{itemize}
\item \textbf{Supervisionado} $\times$ \textbf{não supervisionado}
$\times$ \textbf{semissupervisionado.} Supervisionado usa dados
\textbf{rotulados} (classificação/regressão); não supervisionado acha
estrutura \textbf{sem rótulo} (clusterização/associação). A FGV inverte
``rotulado'' entre os dois.
\item \textbf{Overfitting} (decora o treino, vai mal em dados novos) $\times$
\textbf{underfitting} (modelo simples demais) --- mecanismo detalhado no
Capítulo \ref{cap:atualidades}, seção de viés $\times$ variância.
\item \textbf{Métricas:} MAE = erro médio absoluto (regressão); F1 = média
harmônica de precisão e recall; ROC/AUC = capacidade de separar classes
variando o limiar.
\item \textbf{Redes neurais} (ativação sigmoide) e CNN; \textbf{LLM/IA
generativa:} RAG, agentes, engenharia de prompt, ética e explicabilidade ---
ver também Capítulo \ref{cap:atualidades}.
\item Em cenário de IA em produção, ``monitorar drift'' e deploy
blue-green/canário são as respostas típicas de \textbf{MLOps}. Não confundir
IA generativa (cria conteúdo) com IA discriminativa (classifica/prediz).
\end{itemize}

\begin{conceito}[Regularização --- por que a penalidade encolhe o modelo]
O \emph{overfitting} tem uma assinatura numérica: para passar exatamente por
cada ponto do treino, o modelo precisa de \textbf{coeficientes grandes} --- é
o que permite a curva subir e descer bruscamente entre um ponto e outro.
Regularizar é atacar essa assinatura: em vez de minimizar só o erro, o
treinamento passa a minimizar \textbf{erro + penalidade sobre o tamanho dos
coeficientes}. Coeficiente grande fica caro, e o modelo prefere uma curva mais
lisa --- que generaliza melhor.

O que muda entre as duas famílias é \textbf{como} a penalidade é calculada, e
daí vem a diferença de efeito:

\begin{center}
\begin{tabular}{@{}p{2.4cm}p{3.6cm}p{5.4cm}@{}}
\toprule
& \textbf{Penaliza} & \textbf{Efeito} \\
\midrule
\textbf{L1 (Lasso)} & a soma dos \textbf{valores absolutos} dos coeficientes &
empurra coeficientes \textbf{até zero} --- o atributo sai do modelo. Por isso
L1 faz \textbf{seleção de atributos} e produz modelo \emph{esparso} \\
\textbf{L2 (Ridge)} & a soma dos \textbf{quadrados} dos coeficientes &
\textbf{encolhe} todos proporcionalmente, mas \textbf{não zera} --- todos os
atributos continuam no modelo, com peso menor \\
\bottomrule
\end{tabular}
\end{center}

\textbf{A intuição do porquê:} ao elevar ao quadrado, a L2 pune muito o
coeficiente grande e quase nada o que já está perto de zero --- então não vale
a pena terminar de zerá-lo. A L1 pune o mesmo tanto por unidade em qualquer
faixa, então zerar de vez compensa. É essa diferença de \emph{formato} da
penalidade, e não uma escolha arbitrária, que gera a esparsidade.

Consequência que decide o item: se o cenário pede \textbf{descartar atributos
irrelevantes} de um conjunto com centenas de variáveis, é \textbf{L1}. Se pede
apenas \textbf{estabilizar} um modelo com atributos correlacionados, sem
perder nenhum, é \textbf{L2}.
\end{conceito}

\begin{conceito}[Drift --- o modelo não ``estraga'', o mundo é que anda]
Um modelo em produção perde acurácia com o tempo sem que uma linha de código
mude. A causa é que ele aprendeu uma fotografia e o mundo continuou andando ---
e existem \textbf{duas} coisas diferentes que podem ter andado:

\begin{center}
\begin{tabular}{@{}p{2.6cm}p{4.2cm}p{4.8cm}@{}}
\toprule
& \textbf{O que mudou} & \textbf{Exemplo} \\
\midrule
\textbf{Data drift} & a \textbf{distribuição da entrada} $P(X)$ --- os dados
que chegam já não se parecem com os de treino & o público do sistema mudou:
antes vinham pedidos de servidores de 40--50 anos, agora chegam de 20--30. A
regra ``quem é assim tende a cancelar'' continua valendo \\
\addlinespace
\textbf{Concept drift} & a \textbf{relação entrada$\to$saída} $P(Y|X)$ --- o
próprio conceito-alvo & a entrada é a mesma, mas o comportamento mudou: depois
de uma lei nova, o mesmo perfil de cliente que antes cancelava passou a
permanecer. O modelo aprendeu uma regra que \textbf{deixou de ser verdadeira}
\\
\bottomrule
\end{tabular}
\end{center}

\textbf{Como distinguir na leitura do enunciado:} pergunte \emph{o que o
cenário diz que mudou}. Se o texto descreve \textbf{quem} está entrando no
sistema (perfil, faixa, origem, volume), é \textbf{data} drift. Se descreve
que \textbf{o mesmo tipo de caso passou a ter desfecho diferente} --- nova
legislação, nova campanha do concorrente, pandemia, fraude que mudou de
tática ---, é \textbf{concept} drift.

A diferença tem consequência prática, e é isso que o item de MLOps cobra: data
drift às vezes se resolve \textbf{rebalanceando ou reamostrando} os dados;
concept drift \textbf{exige retreinar com rótulos novos}, porque a resposta
certa mudou. Detectar é papel do monitoramento contínuo --- e é por isso que
``monitorar drift'' e ``gatilho automático de retreino'' são as respostas
típicas em cenário de modelo degradando em produção.
\end{conceito}

\begin{pegadinha}
Inverter hipervisor: \textbf{Tipo 1} (bare-metal) roda direto no hardware;
\textbf{Tipo 2} roda sobre um SO hospedeiro; \textbf{contêiner} compartilha o
kernel do host (não virtualiza hardware). Troca internet/intranet/extranet:
intranet = interna; extranet = estende a intranet a parceiros externos.
\end{pegadinha}

\begin{jacaiu}
\textbf{Em prova real da FGV: 21 questões, e nenhuma delas é conteúdo do
Perfil 3.} Depois que a classificação foi consertada, este bloco virou o que o
nome sempre prometeu --- o que não tem casa --- e ele se divide em duas
famílias, as duas fora do edital:

\begin{itemize}
\item \textbf{Arquitetura de computadores e sistema operacional (11).}
Gargalo de \textbf{Von Neumann}, ULA e unidade de controle, \emph{overflow}
$\times$ \emph{carry}, complemento de dois, conversão de base, ROM/firmware,
periférico, ciclo busca-decodificação-execução, \textbf{MMU/TLB} e
\emph{thrashing}, tabela de páginas, \textbf{DMA} --- \textbf{ALERO 2026} e
\textbf{TJ-RJ}. O edital do Perfil 3 fala em arquitetura \emph{de software}
(Capítulo \ref{cap:arquitetura}).
\item \textbf{Administração e direito público (10).} Avaliação de desempenho
de gestores, priorização de melhoria de processo, liderança de equipe,
departamentalização, Resolução CNMP, redistribuição de cargo, compromisso da
LINDB, critério de desempate em licitação, responsabilidade da concessionária
e processo disciplinar --- \textbf{MPU}, tudo do Módulo I daquele concurso.
\end{itemize}

\textbf{O que saiu daqui.} Até o fechamento do ITEM 7 este bloco declarava 68
questões e ``maioria esmagadora DBA''. Era sintoma de um defeito do
importador: o rótulo de sub-bloco dos mapas da ALERO usava o \emph{slug}
(\texttt{banco-dados}, \texttt{eng-software}), que a tabela de tags não
reconhecia, e 47 questões caíam aqui por engano. Foram para onde sempre
pertenceram --- \textbf{34 para banco de dados} (Capítulo
\ref{cap:banco-dados}) e \textbf{13 para engenharia de software} (Capítulo
\ref{cap:eng-software}). O \emph{conteúdo} de administração física (ARIES,
B+ Tree, tablespaces, otimizador, backup) continua sendo ensinado neste
capítulo, e o Capítulo \ref{cap:banco-dados} remete para cá.

\textbf{No nosso banco} (previsto pelo edital, ainda não visto na amostra de
provas): blockchain pelo encadeamento de hash; RPA; low-code; os Vs do Big
Data; supervisionado $\times$ não supervisionado. Dois itens que já estiveram
nesta lista \textbf{caíram} de fato, só que sob outra tag na
\textbf{Dataprev 2024}: \textbf{servidor web $\times$ servidor de aplicação}
(Capítulo \ref{cap:arquitetura}) e \textbf{internet/intranet/extranet/portal}
(item de redes, Capítulo \ref{cap:redes}).

Regra prática: o que este capítulo \emph{ensina} continua valendo muito ---
administração física de banco de dados é matéria de prova, e as questões dela
agora estão sob \texttt{banco-dados}. O que ficou com a etiqueta de órfã é
material de outro perfil: reconheça e siga em frente. Memorize os pares:
AD $\times$ DBA, incremental $\times$ diferencial,
RPO $\times$ RTO, ordem do ARIES, bitmap $\times$ B+ Tree, rotulado $\times$
não rotulado, data $\times$ concept drift, L1 $\times$ L2, Tipo 1 $\times$
Tipo 2, over $\times$ underfitting.

Rode \texttt{./quiz.py orfaos}.
\end{jacaiu}

\begin{comosair}
Monte flashcards dos pares que a FGV inverte neste bloco (listados acima) e
revise-os isolados dos outros capítulos --- é o jeito mais rápido de fixar um
bloco que mistura tantos assuntos diferentes.
\end{comosair}
